\documentclass{../yalumba}

\title{Coalition Transfer for\\Reference-Free Relatedness Detection}
\author{Ajax Davis}
\date{April 2026}

\setabstract{%
We present \emph{Coalition Transfer}, a reference-free method for detecting
biological relatedness directly from raw sequencing reads.  The algorithm
extends the symbiogenesis family of approaches by shifting the unit of analysis
from individual k-mer modules to \emph{coalitions}---sets of module identifiers
that co-occur within a single sequencing read and are fingerprinted as composite
inheritance units.  Where Module Persistence~v2 tracks which modules appear in
each sample, Coalition Transfer asks a deeper question: which \emph{combinations}
of modules travel together through the inheritance process, preserving their
internal composition across generations?  The method proceeds in four stages:
(1)~module extraction using 150\,bp windows with a minimum support of~5 and
cohesion threshold of~0.25, discovering 507--865 modules per sample from a
shared vocabulary of 3{,}656 unique modules; (2)~coalition profiling, which
fingerprints each read's module set and tracks coalition frequency and integrity
across the sample; (3)~informative coalition filtering, retaining 15{,}181 of
17{,}882 total coalitions that appear in a discriminative subset of samples; and
(4)~a three-component scoring function combining informative Jaccard overlap
(weight~0.25), IDF-weighted rarity (weight~0.45), and integrity cosine
similarity (weight~0.30).  Evaluated on six members of the CEPH~1463
three-generation pedigree (10~pairwise comparisons), Coalition Transfer achieves
an average score of~$0.4064$ for parent--child pairs versus~$0.3928$ for
unrelated and in-law pairs, yielding a mean separation of~$+1.36$ percentage
points.  However, the weakest gap is~$-0.55$ percentage points: the in-law pair
Mat.GF~$\leftrightarrow$~Father ($0.4086$) scores above the weakest
parent--child pair Pat.GM~$\to$~Father ($0.4032$), precluding threshold-based
classification.  Compared to Module Persistence~v2 ($+4.76\%$ separation,
1{,}085\,s prep), Coalition Transfer achieves a~$4{\times}$ speedup
(272\,s total) but lower separation.  Compared to Rare-Run~P90 ($+583\%$
separation), both symbiogenesis methods remain far behind.  We analyze why the
integrity component fails to discriminate (all pairs score~1.0), why
rarity-weighted scoring carries the real signal, and how the persistent in-law
confound limits all current approaches.  The work establishes coalitions as a
meaningful unit of genomic analysis while honestly documenting the gap between
ecological metaphor and practical classification.}

\begin{document}
\maketitle

\tableofcontents
\newpage

% ════════════════════════════════════════════════════════════════════
\section{Introduction}
\label{sec:intro}

% ── The inheritance unit problem ──

\subsection{The Inheritance Unit Problem}

Biological inheritance operates on blocks of DNA, not individual nucleotides.
When a parent transmits a chromosome to a child, an entire haplotype
block---spanning thousands of bases and containing dozens to hundreds of
variant sites---travels as a single unit.  Classical identity-by-descent (IBD)
detection exploits this fact by searching for long shared haplotype segments
between individuals~\cite{browning2012ibd, gusev2009germline}.  The longer the
shared segment, the more recent the common ancestor, and the closer the
biological relationship.

\ymarginnote{Inheritance operates on blocks, not individual bases.}

In reference-free genomics, we have no coordinate system, no phased haplotypes,
and no variant calls.  We operate directly on raw sequencing reads---short
fragments of 100--300 bases each, drawn from unknown positions across the
genome.  The challenge is to reconstruct some analogue of ``shared segments''
from this unordered, unlocalized data.

Our previous work on Module Persistence~\cite{margulis1998} introduced the
first symbiogenesis algorithm: k-mer modules, defined as connected components
of co-occurring k-mers within genomic windows, served as the unit of analysis.
Module sharing between samples provided a meaningful signal for relatedness
detection, achieving $+4.76\%$ separation between parent--child and unrelated
pairs on the CEPH~1463 pedigree.  But modules are still relatively fine-grained
units---each module captures a single cluster of co-occurring k-mers,
analogous to a single ``species'' in the symbiogenesis metaphor.

\subsection{Coalition Transfer: The Second Symbiogenesis Algorithm}

This report introduces \emph{Coalition Transfer}, the second algorithm in the
symbiogenesis family.  The key conceptual advance is a shift from individual
modules to \emph{coalitions}---composite units formed by the set of all modules
that co-occur within a single sequencing read.

\begin{definition}
A \textbf{coalition} is the set of module identifiers extracted from a single
sequencing read.  It represents a composite inheritance unit: a combination of
k-mer modules that were physically proximate in the donor genome and were
captured together on a single read fragment.
\end{definition}

The biological intuition is straightforward.  If a parent transmits a haplotype
block to a child, and that block spans multiple module boundaries, then the
\emph{combination} of modules on reads from that block will be conserved in the
child's data.  A coalition captures exactly this combination.  Two individuals
who share a coalition share not just individual modules (which might appear
anywhere in the genome) but a specific \emph{arrangement} of modules that was
inherited as a unit.

\ymarginnote[yalumba-violet]{Coalitions capture \emph{combinations} of modules, not individual modules.}

The name ``coalition'' is borrowed from ecology and game theory, where it
describes a temporary alliance of distinct entities that cooperate for mutual
benefit.  In our context, the ``entities'' are k-mer modules, and the
``cooperation'' is their physical co-location on inherited DNA.  The coalition
persists across generations---it is ``transferred'' from parent to child---as
long as recombination does not break the underlying haplotype block.

\subsection{What Coalition Transfer Addresses}

Module Persistence~v2 achieved meaningful separation on the CEPH~1463 pedigree,
but several limitations motivated the development of Coalition Transfer:

\begin{enumerate}
  \item \textbf{Module independence assumption.}  Module Persistence treats each
        module as an independent feature.  In reality, modules that co-occur on
        the same read carry joint information that is lost when they are scored
        independently.
  \item \textbf{Preparation time.}  Module Persistence required approximately
        1{,}085 seconds of preparation time (module extraction plus graph
        construction).  A faster pipeline would enable iterative experimentation
        and, eventually, clinical use.
  \item \textbf{Spectral similarity.}  The spectral cosine component in
        Module Persistence requires eigendecomposition of the module co-occurrence
        graph, which is computationally expensive and conceptually opaque.
        Coalition Transfer replaces this with a simpler integrity-based measure.
  \item \textbf{In-law confound.}  The most persistent failure in all yalumba
        experiments is the in-law pair Mat.GF~$\leftrightarrow$~Father scoring
        above some parent--child pairs.  Coalition Transfer explores whether
        composite units provide additional discriminative power.
\end{enumerate}

\subsection{Contributions}

This report makes the following contributions:

\begin{enumerate}
  \item A complete description of the Coalition Transfer algorithm, including
        module extraction with 150\,bp windows, coalition fingerprinting,
        coalition profiling, informative filtering, and three-component scoring.
  \item An evaluation on all 10 pairwise comparisons in the CEPH~1463 pedigree,
        with full component breakdowns and ranking analysis.
  \item A direct comparison to Module Persistence~v2 and Rare-Run P90, with
        honest analysis of where Coalition Transfer improves and where it falls
        short.
  \item An analysis of why the integrity component fails to discriminate, why
        rarity carries the signal, and what this implies for the symbiogenesis
        research program.
  \item Detailed computational profiling showing a $4{\times}$ speedup over
        Module Persistence~v2.
\end{enumerate}

\subsection{Relationship to Prior Experiments}

Coalition Transfer was developed as the logical successor to Module
Persistence~v2 (experiment~50 in the yalumba lab notebook).  The module
extraction phase shares the same algorithmic core but uses different
parameters (150\,bp windows instead of 100\,bp, higher minimum support).  The
coalition profiling and scoring phases are entirely new.  All experiments use
the same six CEPH~1463 samples and the same 10-pair evaluation protocol,
ensuring strict comparability across the entire experimental trajectory from
experiment~1 through the present work.


% ════════════════════════════════════════════════════════════════════
\section{Background}
\label{sec:background}

\subsection{Symbiogenesis Theory}

\ymarginnote{Margulis: complexity arises from merger, not mutation.}

Lynn Margulis's theory of symbiogenesis~\cite{margulis1998} proposed that
the major transitions in biological complexity---the origin of eukaryotic
cells, the acquisition of mitochondria and chloroplasts---arose not through
gradual accumulation of point mutations but through the \emph{merger} of
previously independent organisms into a new, more capable whole.  The
endosymbiotic origin of mitochondria is now established fact: an
$\alpha$-proteobacterium was engulfed by an archaeal host cell, and the
resulting chimera became the ancestor of all eukaryotes.

The symbiogenesis framework provides a powerful metaphor for our algorithmic
approach.  Just as mitochondria are ``organisms within organisms'' that
maintain their own genome while contributing to the host, k-mer modules are
``structures within reads'' that maintain their own identity while contributing
to the overall relatedness signal.  Coalitions extend the metaphor further:
they represent the \emph{cellular organization}---the specific arrangement of
organelles within a cell---that is inherited as a unit across generations.

\subsection{Coalitions in Ecology}

In community ecology, a coalition describes a group of species that co-occur
more frequently than expected by chance~\cite{bray1957, morisita1959}.
Coalitions are distinguished from communities by their \emph{compositionality}:
a coalition is defined not by the presence of any single species but by the
joint presence of a specific combination.  The classic example is the
mycorrhizal coalition---a fungus, a host plant, and specific bacterial
associates---that functions as an integrated unit despite comprising distinct
organisms.

We adopt this ecological definition directly.  Our coalitions are sets of
module IDs that co-occur on a single read, and the ``ecosystem'' is the
genome of an individual.  The transfer of coalitions between parent and child
mirrors the vertical transmission of symbiotic communities in biology, where
offspring inherit not just individual symbionts but entire symbiotic
assemblages.

\subsection{IBD Segments as Genomic Coalitions}

\ymarginnote[yalumba-violet]{IBD segments are, in essence, coalitions of variants.}

Classical IBD segments~\cite{browning2012ibd, browning2013refinedibd} can be
reinterpreted through the coalition lens.  An IBD segment shared between two
individuals is a chromosomal region where the combination of alleles is
identical because both copies descend from a single ancestral chromosome.
This is precisely a coalition: a set of genetic features (alleles at variant
sites) that travel together as a unit through the inheritance process.

The difference between classical IBD and our approach is the feature set.
Classical IBD operates on phased genotypes at known variant positions,
requiring a reference genome and a variant-calling pipeline.  Our coalitions
operate on k-mer modules extracted from raw reads, requiring nothing beyond
the FASTQ files themselves.  The information loss is substantial---we have no
coordinates, no phasing, and no explicit variant calls---but the
\emph{conceptual structure} is identical: we are searching for combinations
of features that are shared because they were inherited together.

\subsection{Why Coalition Integrity Matters}

Not all coalitions are created equal.  A coalition that appears with exactly
the same module composition every time it is observed carries stronger
evidence of inheritance than one whose composition varies across reads.
We formalize this intuition through \emph{coalition integrity}: the fraction
of observations where the coalition's module set is identical to its canonical
form.

High-integrity coalitions are likely to represent genuine haplotype blocks
where the module arrangement is stable across the genome.  Low-integrity
coalitions may arise from sequencing errors, chimeric reads, or regions where
module boundaries are ambiguous.  By weighting coalitions by their integrity,
we aim to up-weight the reliable inheritance signal and down-weight noise.

As we will see in Section~\ref{sec:results}, this intuition does not survive
contact with real data: all coalitions achieve integrity~1.0 in our current
implementation, rendering the integrity component non-discriminative.
Understanding \emph{why} this happens, and how to fix it, is a central
theme of the Discussion.

\subsection{Reference-Free Approaches}

Several groups have explored reference-free methods for genomic comparison.
Mash~\cite{ondov2016mash} uses MinHash sketches to estimate Jaccard
similarity between genomes, providing fast distance estimation for phylogenetic
placement.  Br\v{i}nda et al.~\cite{brinda2023phylogenetic} demonstrated
reference-free phylogenetic surveillance using k-mer-based approaches.
Fan et al.~\cite{fan2021reference} explored alignment-free relatedness
estimation, though their approach still relied on reference-derived features
for calibration.

Our work differs from all of these in its focus on \emph{structural}
features---modules and coalitions---rather than individual k-mer frequencies.
We hypothesize that the combinatorial structure of inherited DNA carries
information that is lost when genomes are reduced to bags of independent
k-mers.  Coalition Transfer is an explicit test of this hypothesis.


% ════════════════════════════════════════════════════════════════════
\section{Methods}
\label{sec:methods}

\subsection{Dataset}
\label{sec:methods-dataset}

We evaluate on six members of the CEPH~1463 pedigree~\cite{ceph1463},
a three-generation family that has been extensively studied by the Genome
in a Bottle (GIAB) consortium~\cite{zook2019giab} and the 1000 Genomes
Project~\cite{1000genomes2015}.  The six individuals comprise two sets of
grandparents, the father, and the mother:

\begin{pedigree}
\begin{verbatim}
NA12889 + NA12890          NA12891 + NA12892
(Pat.GF)  (Pat.GM)         (Mat.GF)  (Mat.GM)
        |                          |
     NA12877 (Father) ——— NA12878 (Mother)
\end{verbatim}
\end{pedigree}

\noindent This pedigree yields 10 distinct pairwise comparisons:

\begin{itemize}
  \item \textbf{4 parent--child pairs} (marked with $\blacktriangleright$):
    Pat.GF~$\to$~Father, Pat.GM~$\to$~Father,
    Mat.GF~$\to$~Mother, Mat.GM~$\to$~Mother.
  \item \textbf{2 in-law pairs:}
    Mat.GF~$\leftrightarrow$~Father, Pat.GF~$\leftrightarrow$~Mother.
  \item \textbf{3 cross-family pairs:}
    Pat.GF~$\leftrightarrow$~Mat.GF,
    Pat.GF~$\leftrightarrow$~Mat.GM,
    Pat.GM~$\leftrightarrow$~Mat.GM.
  \item \textbf{1 spouse pair:}
    Father~$\leftrightarrow$~Mother (unrelated).
\end{itemize}

All samples are whole-genome sequencing at approximately 50$\times$ coverage,
sequenced on Illumina platforms with paired-end 150\,bp reads.

\subsection{Module Extraction}
\label{sec:methods-modules}

Module extraction follows the same conceptual pipeline as Module
Persistence~v2 but with modified parameters optimized for Coalition Transfer:

\begin{enumerate}
  \item \textbf{K-mer extraction.}  For each read, we extract all canonical
        15-mers using a Rabin--Karp rolling hash~\cite{karp1987rabin} with
        base prime~31 and Mersenne prime $2^{61}-1$.  Canonical k-mers are
        defined as the lexicographically smaller of a k-mer and its reverse
        complement.
  \item \textbf{Frequency filtering.}  We retain only k-mers that appear in
        at least \texttt{minSupport}~$= 5$ reads, removing sequencing errors
        and ultra-rare variants.
  \item \textbf{Co-occurrence counting.}  Within a sliding window of
        \texttt{windowSize}~$= 150$\,bp, we count pairwise co-occurrences of
        frequent k-mers, accumulating counts in a fixed-size hash table.
  \item \textbf{Module discovery.}  K-mer pairs with co-occurrence rates
        exceeding \texttt{minCohesion}~$= 0.25$ (i.e., they appear together in
        at least 25\% of the windows where either appears) are connected by
        edges, and modules are defined as connected components of this graph.
\end{enumerate}

\ymarginnote{150\,bp windows capture the full read length, maximizing co-occurrence signal.}

The key parameter change from Module Persistence~v2 is the window size:
150\,bp instead of 100\,bp.  This matches the typical Illumina read length,
ensuring that co-occurrence statistics reflect the full span of each read.
The larger window increases the number of k-mer pairs considered per window
but reduces the total number of windows, resulting in faster overall execution.

\begin{definition}
A \textbf{module} $M$ is a connected component of the k-mer co-occurrence
graph $G = (V, E)$, where $V$ is the set of frequent k-mers and
$(u, v) \in E$ if and only if
$\mathrm{cohesion}(u, v) \geq \texttt{minCohesion}$.
The cohesion is defined as:
\begin{equation}
  \mathrm{cohesion}(u, v) = \frac{\mathrm{cooccur}(u, v)}{\max(\mathrm{freq}(u), \mathrm{freq}(v))}
  \label{eq:cohesion}
\end{equation}
\end{definition}

\subsection{Shared Vocabulary Construction}
\label{sec:methods-vocabulary}

After modules are extracted independently for each sample, we construct a
shared vocabulary by assigning canonical identifiers to modules that are
structurally equivalent across samples.  Two modules are considered equivalent
if they share a sufficient proportion of their constituent k-mers (measured
by Jaccard overlap of the k-mer sets).

The shared vocabulary for the CEPH~1463 dataset contains 3{,}656 unique
modules.  This shared vocabulary serves as the ``alphabet'' for coalition
construction: only modules present in the shared vocabulary are used to
define coalitions, ensuring that pairwise comparisons operate over a common
feature space.

\subsection{Coalition Definition}
\label{sec:methods-coalition}

\begin{definition}
Given a sequencing read $r$ and a set of modules $\mathcal{M}$ from the
shared vocabulary, the \textbf{coalition} of $r$ is the set of module IDs
whose constituent k-mers appear in $r$:
\begin{equation}
  C(r) = \{m \in \mathcal{M} \mid \exists\, \text{k-mer } k \in m \text{ such that } k \in r\}
  \label{eq:coalition-def}
\end{equation}
Two reads have the same coalition if and only if they produce identical
module-ID sets.
\end{definition}

Coalition identity is established by fingerprinting: the sorted set of module
IDs is hashed to produce a single 64-bit identifier.  This allows efficient
lookup and counting without storing the full module set for each read.

The number of possible coalitions is combinatorial in the vocabulary size
(up to $2^{3656}$ in theory), but in practice the number of observed
coalitions is much smaller, bounded by the number of reads in the dataset.
For our samples, we observe 17{,}882 total distinct coalitions across all
six individuals.

\subsection{Coalition Profiling}
\label{sec:methods-profiling}

For each sample, we construct a \emph{coalition profile} that records:

\begin{enumerate}
  \item \textbf{Coalition frequency:} the number of reads assigned to each
        coalition fingerprint.
  \item \textbf{Coalition integrity:} the fraction of reads assigned to a
        coalition whose module set exactly matches the canonical form.
\end{enumerate}

The integrity measure is intended to capture how ``clean'' a coalition is.
A coalition with integrity~1.0 means that every read assigned to it produced
exactly the same module set, with no partial matches or edge effects.  In our
current implementation, coalition identity is determined by exact module-set
match (via fingerprinting), so integrity is mechanically~1.0 for all
coalitions.  We retain the measure as a placeholder for future refinements
where approximate matching may produce non-trivial integrity values.

Coalition profiling runs in approximately 1.1 seconds per sample (6.6 seconds
total for six samples), making it the fastest stage of the pipeline.

\subsection{Informative Coalition Filtering}
\label{sec:methods-filtering}

Not all coalitions carry relatedness signal.  Coalitions that appear in all
samples reflect species-wide genomic structure (e.g., conserved regulatory
regions) rather than inherited variation.  Coalitions that appear in only one
sample may reflect sequencing artifacts or private mutations that are not
informative for relatedness.

We define a coalition as \textbf{informative} if it appears in at least~2
but not all~6 samples.  This is the same filtering criterion used for
informative modules in Module Persistence~v2, applied at the coalition level.
Of 17{,}882 total coalitions, 15{,}181 (84.9\%) pass the informativeness
filter.

\begin{equation}
  \mathcal{C}_{\mathrm{info}} = \{c \in \mathcal{C} \mid 2 \leq |\{s : c \in s\}| < N\}
  \label{eq:informative}
\end{equation}

\noindent where $N = 6$ is the total number of samples.

\subsection{Three-Component Scoring}
\label{sec:methods-scoring}

The pairwise relatedness score between samples $A$ and $B$ combines three
components:

\begin{equation}
  S(A, B) = 0.25 \cdot J_{\mathrm{info}} + 0.45 \cdot R + 0.30 \cdot I_{\cos}
  \label{eq:score}
\end{equation}

\subsubsection{Component 1: Informative Jaccard ($J_{\mathrm{info}}$)}

The Jaccard index over informative coalitions shared between two samples:

\begin{equation}
  J_{\mathrm{info}}(A, B) = \frac{|\mathcal{C}_{\mathrm{info}}^A \cap \mathcal{C}_{\mathrm{info}}^B|}{|\mathcal{C}_{\mathrm{info}}^A \cup \mathcal{C}_{\mathrm{info}}^B|}
  \label{eq:info-jaccard}
\end{equation}

This measures the raw overlap in coalition vocabularies, without regard to
how rare or common each coalition is.  It receives the lowest weight (0.25)
because, like module-level Jaccard, it is dominated by moderately common
coalitions that provide limited discriminative power.

\subsubsection{Component 2: IDF-Weighted Rarity Score ($R$)}

To up-weight coalitions that are rare across the sample set, we apply
inverse document frequency (IDF) weighting~\cite{salton1988tfidf}:

\begin{equation}
  \mathrm{IDF}(c) = \log\!\left(\frac{N}{|\{s : c \in s\}|}\right)
  \label{eq:idf}
\end{equation}

\begin{equation}
  R(A, B) = \frac{\sum_{c \in \mathcal{C}_{\mathrm{info}}^A \cap \mathcal{C}_{\mathrm{info}}^B} \mathrm{IDF}(c)}{\sum_{c \in \mathcal{C}_{\mathrm{info}}^A \cup \mathcal{C}_{\mathrm{info}}^B} \mathrm{IDF}(c)}
  \label{eq:rarity}
\end{equation}

The rarity score receives the highest weight (0.45) based on the hypothesis
that shared rare coalitions carry stronger evidence of recent common ancestry
than shared common coalitions.  A coalition that appears in only 2 of 6
samples and is shared between a parent--child pair provides strong evidence
that it was inherited, whereas a coalition appearing in 5 of 6 samples is
likely a common genomic feature.

\subsubsection{Component 3: Integrity Cosine ($I_{\cos}$)}

For coalitions shared between $A$ and $B$, we compute the cosine similarity
of their integrity vectors:

\begin{equation}
  I_{\cos}(A, B) = \frac{\mathbf{i}_A \cdot \mathbf{i}_B}{\|\mathbf{i}_A\| \cdot \|\mathbf{i}_B\|}
  \label{eq:integrity-cosine}
\end{equation}

\noindent where $\mathbf{i}_A$ and $\mathbf{i}_B$ are vectors of integrity
values for all shared coalitions.  This component is intended to measure
whether shared coalitions have similar ``quality'' in both samples.  As noted
in Section~\ref{sec:methods-profiling}, all integrity values are~1.0 in the
current implementation, so $I_{\cos} = 1.0$ for all pairs.  The component
contributes a constant $0.30 \times 1.0 = 0.30$ to every score.

\subsection{Evaluation Metrics}
\label{sec:methods-eval}

We evaluate the scoring function using two metrics:

\begin{enumerate}
  \item \textbf{Mean separation:} the difference between the average score
        of the four parent--child pairs and the average score of the six
        unrelated/in-law pairs:
        \begin{equation}
          \Delta_{\mathrm{avg}} = \overline{S}_{\mathrm{related}} - \overline{S}_{\mathrm{unrelated}}
          \label{eq:separation}
        \end{equation}
  \item \textbf{Weakest gap:} the difference between the lowest-scoring
        parent--child pair and the highest-scoring unrelated/in-law pair:
        \begin{equation}
          \Delta_{\mathrm{weak}} = \min(S_{\mathrm{related}}) - \max(S_{\mathrm{unrelated}})
          \label{eq:weakest-gap}
        \end{equation}
        A negative weakest gap means that at least one unrelated pair scores
        above at least one parent--child pair, precluding threshold-based
        classification.
\end{enumerate}

\subsection{Implementation}
\label{sec:methods-impl}

The entire pipeline is implemented in TypeScript running on the Bun runtime,
consistent with the yalumba project's commitment to a vertically integrated
stack with no external bioinformatics libraries.

Key implementation details:

\begin{itemize}
  \item \textbf{Hash table:} Co-occurrence counts are stored in a fixed-size
        hash table of 16{,}777{,}216 entries ($2^{24}$), using open addressing
        with linear probing.  At peak fill, approximately 11.7 million entries
        are occupied (70\% load factor).
  \item \textbf{Memory management:} Module extraction uses typed arrays
        (\texttt{Uint32Array}, \texttt{Float64Array}) throughout, minimizing
        garbage collection pressure.  The pipeline achieves approximately 0\,MB
        heap delta across the full run, indicating efficient allocation and
        deallocation.
  \item \textbf{Coalition fingerprinting:} Module ID sets are sorted and
        hashed using the same Rabin--Karp polynomial hash used for k-mers,
        producing a 64-bit fingerprint per coalition.
  \item \textbf{Streaming architecture:} Reads are processed in a single pass
        per sample.  Module extraction streams through the FASTQ file,
        coalition profiling operates on the extracted module set, and
        comparison iterates over the profile dictionaries.
\end{itemize}


% ════════════════════════════════════════════════════════════════════
\section{Results}
\label{sec:results}

\subsection{Module Extraction Statistics}
\label{sec:results-modules}

Table~\ref{tab:module-stats} summarizes the module extraction results for
each of the six CEPH~1463 samples.

\begin{table}[htbp]
\centering
\tablestyle
\caption{Module extraction statistics per sample.  Window size = 150\,bp,
motif $k$ = 15, minimum support = 5, minimum cohesion = 0.25.}
\label{tab:module-stats}
\begin{tabular}{lrrrr}
\toprule
\rowcolor{table-header}
\textcolor{white}{\textbf{Sample}} &
\textcolor{white}{\textbf{Freq.\ motifs}} &
\textcolor{white}{\textbf{Total motifs}} &
\textcolor{white}{\textbf{Pair entries}} &
\textcolor{white}{\textbf{Modules}} \\
\midrule
NA12889 (Pat.GF)  & $\sim$215k & $\sim$500k & $\sim$11.7M & 627 \\
NA12890 (Pat.GM)  & $\sim$173k & $\sim$455k & $\sim$11.7M & 507 \\
NA12891 (Mat.GF)  & $\sim$232k & $\sim$519k & $\sim$11.7M & 712 \\
NA12892 (Mat.GM)  & $\sim$259k & $\sim$539k & $\sim$11.7M & 865 \\
NA12877 (Father)  & $\sim$198k & $\sim$480k & $\sim$11.7M & 589 \\
NA12878 (Mother)  & $\sim$221k & $\sim$510k & $\sim$11.7M & 695 \\
\bottomrule
\end{tabular}
\end{table}

The number of discovered modules varies considerably across samples
(507--865), reflecting differences in coverage depth, read quality, and
individual genomic complexity.  The pair table consistently fills to
approximately 11.7 million entries (70\% of the 16M-entry table), indicating
that the hash table size is well-calibrated for this dataset.

\subsection{Coalition Profile Statistics}
\label{sec:results-profiles}

Coalition profiling across all six samples yields 17{,}882 total distinct
coalitions, of which 15{,}181 (84.9\%) are classified as informative (appearing
in 2--5 samples).  The remaining 2{,}701 coalitions either appear in all six
samples (universally shared, non-informative) or in only one sample (private,
non-informative).

The shared vocabulary underlying these coalitions comprises 3{,}656 unique
modules---the common language of k-mer structures across the pedigree.  The
fact that $84.9\%$ of coalitions are informative indicates that the coalition
level of analysis naturally selects for discriminative features: most
module \emph{combinations} are not universal, even when the individual modules
might be.

\subsection{Full Score Table}
\label{sec:results-scores}

Table~\ref{tab:scores} presents the complete results for all 10 pairwise
comparisons, ranked by composite score.  Parent--child pairs are marked with
$\blacktriangleright$.

\begin{table}[htbp]
\centering
\tablestyle
\caption{Coalition Transfer scores for all 10 pairwise comparisons in the
CEPH~1463 pedigree, ranked by composite score $S$.  Components:
$J_\mathrm{info}$ = informative Jaccard fraction,
$R$ = IDF-weighted rarity score,
$I_{\mathrm{cos}}$ = integrity cosine (all 1.0).
Parent--child pairs marked $\blacktriangleright$.}
\label{tab:scores}
\begin{tabular}{clcccc}
\toprule
\rowcolor{table-header}
\textcolor{white}{\textbf{Rank}} &
\textcolor{white}{\textbf{Pair}} &
\textcolor{white}{\textbf{$J_\mathrm{info}$}} &
\textcolor{white}{\textbf{$R$}} &
\textcolor{white}{\textbf{$I_{\mathrm{cos}}$}} &
\textcolor{white}{\textbf{$S$}} \\
\midrule
$\blacktriangleright$ 1 & Pat.GF $\to$ Father     & 25.0\% & 10.7\% & 1.0000 & 0.4107 \\
2                       & Mat.GF $\leftrightarrow$ Father (in-law)  & 26.0\% & 9.7\% & 1.0000 & 0.4086 \\
$\blacktriangleright$ 3 & Mat.GM $\to$ Mother      & 24.9\% & 9.7\% & 1.0000 & 0.4061 \\
$\blacktriangleright$ 4 & Mat.GF $\to$ Mother      & 24.9\% & 9.6\% & 1.0000 & 0.4054 \\
$\blacktriangleright$ 5 & Pat.GM $\to$ Father      & 24.6\% & 9.2\% & 1.0000 & 0.4032 \\
6                       & Spouses (unrelated)       & 23.8\% & 8.9\% & 1.0000 & 0.3996 \\
7                       & Pat.GF $\leftrightarrow$ Mat.GM           & 22.7\% & 9.4\% & 1.0000 & 0.3993 \\
8                       & Pat.GM $\leftrightarrow$ Mat.GM           & 23.3\% & 8.5\% & 1.0000 & 0.3966 \\
9                       & Pat.GF $\leftrightarrow$ Mat.GF           & 18.9\% & 7.1\% & 1.0000 & 0.3792 \\
10                      & Pat.GF $\leftrightarrow$ Mother (in-law)  & 17.5\% & 6.6\% & 1.0000 & 0.3732 \\
\bottomrule
\end{tabular}
\end{table}

\subsection{Score Ranking Analysis}
\label{sec:results-ranking}

The ranking reveals a partially successful separation.  The top-scoring pair
is the parent--child pair Pat.GF~$\to$~Father ($S = 0.4107$), and three of
the four parent--child pairs occupy ranks 1, 3, and 4.  However, the in-law
pair Mat.GF~$\leftrightarrow$~Father ($S = 0.4086$) intrudes at rank~2,
splitting the parent--child cluster.

The weakest parent--child pair, Pat.GM~$\to$~Father ($S = 0.4032$), sits at
rank~5---just 0.36 percentage points above the highest-scoring unrelated pair
(Spouses at $S = 0.3996$).  This narrow margin, combined with the in-law
intrusion at rank~2, means that no single threshold can perfectly separate
related from unrelated pairs.

\ymarginnote{Three of four parent--child pairs rank in the top 4, but an in-law intrudes at rank~2.}

Notably, the six unrelated/in-law pairs span a wide range: from 0.3732
(Pat.GF~$\leftrightarrow$~Mother) to 0.4086 (Mat.GF~$\leftrightarrow$~Father),
a spread of 3.54 percentage points.  This within-class variance exceeds the
between-class separation of 1.36 percentage points, confirming that the
current scoring function cannot support binary classification.

\subsection{Separation Analysis}
\label{sec:results-separation}

\begin{keyfinding}
Coalition Transfer achieves a mean separation of $+1.36$ percentage points
between parent--child pairs (average $S = 0.4064$) and unrelated/in-law pairs
(average $S = 0.3928$).  The four parent--child pairs score, on average,
higher than the six unrelated pairs, confirming that the coalition-level
signal captures some genuine inheritance information.
\end{keyfinding}

The informative coalition counts for the top-scoring parent--child pair
(Pat.GF~$\to$~Father) are 2{,}041 shared out of 8{,}152 in the union (25.0\%
informative Jaccard), with a rarity score of 10.7\%.  The bottom-scoring
pair (Pat.GF~$\leftrightarrow$~Mother) shares only 1{,}484 informative
coalitions out of 8{,}492 (17.5\% Jaccard, 6.6\% rarity).  The rarity
component shows the larger dynamic range: 10.7\% vs.\ 6.6\% (4.1~pp spread)
compared to 25.0\% vs.\ 17.5\% (7.5~pp spread) for Jaccard.  However, the
rarity component receives a higher weight (0.45 vs.\ 0.25), amplifying its
contribution to the final score.

\subsection{Weakest Gap Failure}
\label{sec:results-gap}

\begin{limitation}
The weakest gap is $-0.55$ percentage points: the in-law pair
Mat.GF~$\leftrightarrow$~Father ($S = 0.4086$) scores above the weakest
parent--child pair Pat.GM~$\to$~Father ($S = 0.4032$).  This means Coalition
Transfer \textbf{cannot} separate all related pairs from all unrelated pairs
using a single threshold.
\end{limitation}

The problematic in-law pair, Mat.GF~$\leftrightarrow$~Father, has an
informative Jaccard of 26.0\%---actually the highest of any pair in the
dataset.  Its rarity score of 9.7\% exceeds that of two parent--child pairs
(Mat.GF~$\to$~Mother at 9.6\% and Pat.GM~$\to$~Father at 9.2\%).  This
in-law pair shares 1{,}952 informative coalitions out of 7{,}515 in the
union---a substantial overlap that is not explained by direct inheritance.

The persistence of this confound across all yalumba experiments (raw Jaccard,
frequency filtering, run-length scoring, Module Persistence, and now Coalition
Transfer) suggests that Mat.GF and Father share an unusual amount of genomic
structure that is not attributable to their relationship through Mother.
Possible explanations include shared population background (both are of
European ancestry from the Utah CEPH collection), cryptic distant relatedness,
or systematic sequencing artifacts that create spurious module overlap.

\subsection{Component Contribution Analysis}
\label{sec:results-components}

Since the integrity cosine is uniformly~1.0, the effective scoring function
reduces to a two-component formula:

\begin{equation}
  S_{\mathrm{eff}}(A, B) = 0.25 \cdot J_{\mathrm{info}} + 0.45 \cdot R + 0.30
  \label{eq:effective-score}
\end{equation}

The constant $0.30$ term from integrity contributes equally to all pairs and
has no discriminative power.  All separation is driven by the Jaccard and
rarity components, with rarity receiving $1.8{\times}$ the weight of Jaccard.

To understand which component drives the separation, we can decompose the
mean separation:

\begin{itemize}
  \item \textbf{Jaccard contribution:}
    $\overline{J}_{\mathrm{related}} - \overline{J}_{\mathrm{unrelated}}
    \approx 24.85\% - 22.03\% = +2.82\,\mathrm{pp}$,
    weighted: $0.25 \times 2.82 = +0.71\,\mathrm{pp}$.
  \item \textbf{Rarity contribution:}
    $\overline{R}_{\mathrm{related}} - \overline{R}_{\mathrm{unrelated}}
    \approx 9.80\% - 8.37\% = +1.43\,\mathrm{pp}$,
    weighted: $0.45 \times 1.43 = +0.64\,\mathrm{pp}$.
  \item \textbf{Integrity contribution:}
    $1.0 - 1.0 = 0.00$, weighted: $0.30 \times 0.00 = 0.00\,\mathrm{pp}$.
\end{itemize}

The Jaccard and rarity components contribute approximately equally to the
total separation ($+0.71$ vs.\ $+0.64$ pp), despite rarity having a higher
weight.  This suggests that the IDF weighting is not providing as much
additional discriminative power as hoped---the rarity signal is compressed
relative to the raw Jaccard signal.

\subsection{Comparison to Module Persistence v2}
\label{sec:results-mpg}

Table~\ref{tab:comparison-mpg} compares Coalition Transfer to Module
Persistence~v2 on the same dataset and evaluation protocol.

\begin{table}[htbp]
\centering
\tablestyle
\caption{Comparison of Coalition Transfer and Module Persistence~v2.}
\label{tab:comparison-mpg}
\begin{tabular}{lcc}
\toprule
\rowcolor{table-header}
\textcolor{white}{\textbf{Metric}} &
\textcolor{white}{\textbf{Coalition Transfer}} &
\textcolor{white}{\textbf{Module Persistence v2}} \\
\midrule
Mean separation        & $+1.36\%$  & $+4.76\%$ \\
Weakest gap            & $-0.55\%$  & $-3.17\%$ \\
Prep time (total)      & 272\,s     & 1{,}085\,s \\
Prep time (per sample) & $\sim$44\,s & $\sim$180\,s \\
Comparison time        & $<0.01$\,s  & $<0.01$\,s \\
Heap delta             & $\sim$0\,MB & $\sim$4{,}806\,MB \\
Unit of analysis       & Coalition   & Module \\
Components             & 3 (J, R, I) & 3 ($J_M$, $J_E$, $\cos$) \\
\bottomrule
\end{tabular}
\end{table}

Coalition Transfer achieves a $4{\times}$ speedup in preparation time and
dramatically lower memory usage, but at the cost of reduced separation:
$+1.36\%$ vs.\ $+4.76\%$, a factor of $3.5{\times}$ worse.  The weakest
gap, however, is actually less negative ($-0.55\%$ vs.\ $-3.17\%$),
suggesting that Coalition Transfer produces a tighter score distribution
overall---the reduced separation comes from compressing both the related
and unrelated score ranges.

\begin{keyfinding}
Coalition Transfer is $4{\times}$ faster than Module Persistence~v2 and uses
negligible additional memory, but achieves $3.5{\times}$ lower mean
separation.  The weakest gap is less negative ($-0.55\%$ vs.\ $-3.17\%$),
indicating a more compressed but more uniform score distribution.
\end{keyfinding}

\subsection{Comparison to Rare-Run P90 Baseline}
\label{sec:results-rarerun}

Table~\ref{tab:comparison-rarerun} compares Coalition Transfer to the
Rare-Run P90 baseline (experiment~38), which remains the strongest
performer in the yalumba experiment history.

\begin{table}[htbp]
\centering
\tablestyle
\caption{Comparison of Coalition Transfer and Rare-Run P90.}
\label{tab:comparison-rarerun}
\begin{tabular}{lcc}
\toprule
\rowcolor{table-header}
\textcolor{white}{\textbf{Metric}} &
\textcolor{white}{\textbf{Coalition Transfer}} &
\textcolor{white}{\textbf{Rare-Run P90}} \\
\midrule
Mean separation        & $+1.36\%$   & $+583\%$ \\
Weakest gap            & $-0.55\%$   & $+300\%$ \\
Prep time              & 272\,s      & 72.5\,s \\
Unit of analysis       & Coalition   & Run of rare k-mers \\
\bottomrule
\end{tabular}
\end{table}

The performance gap between Coalition Transfer and Rare-Run P90 is enormous:
$+583\%$ vs.\ $+1.36\%$ mean separation, and $+300\%$ vs.\ $-0.55\%$ weakest
gap.  Rare-Run P90 achieves complete separation (positive weakest gap) in
less time.  This comparison underscores a persistent finding across the
symbiogenesis algorithms: the module/coalition abstraction layer does not
yet capture the discriminative signal as effectively as direct rare k-mer
run detection.

\subsection{Timing and Memory Breakdown}
\label{sec:results-timing}

Table~\ref{tab:timing} provides a detailed breakdown of execution time by
pipeline stage.

\begin{table}[htbp]
\centering
\tablestyle
\caption{Execution time breakdown for Coalition Transfer on the CEPH~1463
dataset (6 samples, 10 pairs).}
\label{tab:timing}
\begin{tabular}{lrr}
\toprule
\rowcolor{table-header}
\textcolor{white}{\textbf{Stage}} &
\textcolor{white}{\textbf{Total time}} &
\textcolor{white}{\textbf{Per sample/pair}} \\
\midrule
Data loading          & 24.0\,s   & 4.0\,s per sample \\
Module extraction     & 265.0\,s  & $\sim$44\,s per sample \\
Coalition profiling   & 6.6\,s    & $\sim$1.1\,s per sample \\
Pairwise comparison   & $<0.1$\,s & $<0.01$\,s per pair \\
\midrule
\textbf{Total}        & \textbf{272.1\,s} & --- \\
\bottomrule
\end{tabular}
\end{table}

Module extraction dominates the runtime at 97.4\% of total time.  This is
expected: module extraction involves streaming through every read in every
FASTQ file, extracting k-mers, counting co-occurrences, and building the
module graph.  Coalition profiling, by contrast, operates on the already-extracted
module assignments and is effectively a counting operation.  Pairwise comparison
is instantaneous because it operates on pre-computed coalition profiles.

The heap delta of approximately 0\,MB indicates that the pipeline's memory
management is effective: allocations during module extraction are released
before coalition profiling begins, and the profiling stage uses compact
data structures (hash maps of coalition fingerprints to counts).


% ════════════════════════════════════════════════════════════════════
\section{Discussion}
\label{sec:discussion}

\subsection{What Coalition Transfer Captures}
\label{sec:discussion-captures}

Coalition Transfer captures a genuinely different signal than Module
Persistence.  Where Module Persistence asks ``which modules are present?'',
Coalition Transfer asks ``which \emph{combinations} of modules are present?''.
This is a strictly higher-order question: two samples can share identical
module inventories but have completely different coalition profiles, if the
modules co-occur in different arrangements across reads.

The positive mean separation ($+1.36\%$) confirms that coalition-level
sharing is correlated with biological relatedness.  Parent--child pairs share
more informative coalitions, and the shared coalitions are rarer on average,
than those shared by unrelated pairs.  This is consistent with the inheritance
model: a child inherits specific haplotype blocks from each parent, and the
module combinations within those blocks are transferred as units.

However, the magnitude of the signal is disappointing.  The $+1.36\%$
separation is substantially lower than the $+4.76\%$ achieved by Module
Persistence~v2, which operates at a coarser level of analysis (individual
modules rather than module combinations).  This suggests that the additional
combinatorial specificity of coalitions introduces noise faster than it
introduces signal.

\subsection{Why 150\,bp Windows Improve Prep Time but Not Separation}
\label{sec:discussion-windows}

The switch from 100\,bp to 150\,bp windows was motivated by two goals:
matching the window size to the read length (maximizing co-occurrence
signal per window) and reducing the total number of windows (improving
execution speed).

\ymarginnote{Larger windows trade resolution for speed.}

The speed improvement is substantial: $\sim$44\,s per sample vs.\
$\sim$180\,s, a $4{\times}$ speedup.  But the separation does not improve.
This apparent paradox is explained by the relationship between window size
and module resolution.

With a 150\,bp window, more k-mer pairs fall within the same window, and
more pairs exceed the cohesion threshold.  This produces larger, more
connected modules---fewer modules with more k-mers each.  The modules are
coarser-grained, capturing broader genomic regions but losing the fine
structure that distinguishes related from unrelated sharing.  At the
coalition level, coarser modules mean that coalitions become less specific:
a coalition containing a large, generic module is less informative than one
containing several small, specific modules.

The implication is that window size creates a trade-off between computational
efficiency and signal resolution.  The optimal window size for Coalition
Transfer may be smaller than 150\,bp, at the cost of longer extraction time.

\subsection{The Integrity Measurement Problem}
\label{sec:discussion-integrity}

The most striking failure in Coalition Transfer is the integrity component.
All coalitions achieve integrity~1.0 for all samples, rendering the integrity
cosine a constant and eliminating 30\% of the scoring function's potential
discriminative power.

This is a \emph{measurement artifact}, not a biological finding.  In the
current implementation, coalition identity is determined by exact module-set
match via fingerprinting.  Two reads that differ by even a single module
produce different fingerprints and are assigned to different coalitions.
Consequently, every observation of a given coalition fingerprint has, by
construction, exactly the same module set---integrity is mechanically~1.0.

\begin{limitation}
The integrity component is non-discriminative because coalition identity is
defined by exact fingerprint matching.  Any variation in module composition
creates a new coalition rather than reducing the integrity of an existing one.
This wastes 30\% of the scoring function's weight on a constant.
\end{limitation}

To make integrity meaningful, the implementation would need to support
\emph{approximate coalition matching}: grouping reads whose module sets are
similar but not identical, and measuring what fraction of each coalition's
reads match its canonical form exactly.  This could be achieved through
locality-sensitive hashing (LSH) of module sets, hierarchical clustering
of coalitions, or a Hamming-distance threshold on fingerprints.

\subsection{Rarity-Weighted Scoring as the Real Signal Carrier}
\label{sec:discussion-rarity}

With integrity contributing nothing, the effective scoring function is:

\[
  S_{\mathrm{eff}} = 0.25 \cdot J_{\mathrm{info}} + 0.45 \cdot R + 0.30
\]

The rarity component ($R$) carries 63.4\% of the non-constant weight
($0.45 / (0.25 + 0.45) = 0.643$), making it the dominant signal.  This is
consistent with the broader finding across yalumba experiments that
\emph{rarity is the most informative feature for relatedness detection}.

The IDF-weighted rarity score amplifies the contribution of coalitions that
appear in few samples.  A coalition shared between only a parent--child pair
and no other individuals receives a high IDF weight, contributing strongly
to their pairwise score.  A coalition shared by five of six samples receives
a low IDF weight, contributing weakly regardless of which pair is being
compared.

The connection to information theory is direct: rare events carry more
information~\cite{li2004ncd}.  In the genomic context, a rare coalition is
one that is unlikely to appear by chance in unrelated individuals, so its
shared presence between two individuals is strong evidence of shared ancestry.
This principle is the same one that makes rare variant sharing a powerful
signal in classical IBD detection~\cite{browning2012ibd}.

\subsection{The In-Law Confound Persists}
\label{sec:discussion-inlaw}

The in-law pair Mat.GF~$\leftrightarrow$~Father has confounded every yalumba
experiment since experiment~1.  In Coalition Transfer, this pair scores
$0.4086$---higher than two of four parent--child pairs.  Its informative
Jaccard of 26.0\% is the highest of any pair, suggesting that Mat.GF and
Father share an unusually large number of coalition types.

Several hypotheses could explain this persistent confound:

\begin{enumerate}
  \item \textbf{Shared European ancestry.}  The CEPH collection was drawn
        from Utah families of primarily northern European descent.  Mat.GF
        and Father may share population-level haplotype blocks that create
        genuine coalition overlap without recent common ancestry.
  \item \textbf{Assortative mating.}  If Mother was drawn to a partner with
        similar genomic background to her father, Mat.GF and Father would
        share more genomic structure than expected by chance.
  \item \textbf{Coverage artifacts.}  Differences in sequencing depth between
        samples can inflate or deflate module counts, creating spurious
        overlap between samples with similar coverage profiles.
  \item \textbf{Cryptic relatedness.}  Mat.GF and Father may be distantly
        related through ancestors not captured in the pedigree.
\end{enumerate}

None of these hypotheses can be tested with the current dataset alone.
Resolving the in-law confound likely requires either a larger pedigree (where
the effect can be assessed statistically) or integration with variant-level
information (to distinguish population-level sharing from IBD).

\subsection{Coalition vs.\ Module: Which Unit Is More Informative?}
\label{sec:discussion-unit}

The direct comparison between Coalition Transfer ($+1.36\%$ separation) and
Module Persistence~v2 ($+4.76\%$ separation) strongly favors the module as
the unit of analysis.  This is a surprising finding: we hypothesized that
coalitions, as higher-order structures, would capture more inheritance
information than individual modules.

The explanation likely lies in the \emph{specificity--sensitivity trade-off}.
Coalitions are highly specific: two individuals must share exactly the same
combination of modules (as determined by fingerprinting) for a coalition match
to count.  This specificity means that many genuine inheritance events are
missed because the inherited haplotype block spans slightly different module
boundaries in the two samples, producing different coalitions despite
substantial overlap.

Modules, by contrast, are less specific but more sensitive.  Two individuals
who share a haplotype block will share most of the individual modules within
that block, even if the block boundaries differ slightly.  Module-level
analysis tolerates boundary variation in a way that coalition-level analysis
does not.

\ymarginnote{Coalitions are too specific---exact module-set matching misses genuine inheritance.}

This analysis suggests a potential improvement: \emph{fuzzy coalition
matching}, where two coalitions are considered equivalent if their module sets
overlap above some threshold (e.g., Jaccard $> 0.8$) rather than requiring
exact identity.  This would reduce specificity while maintaining the
compositional structure that makes coalitions conceptually appealing.

\subsection{Connection to Classical Haplotype Block Theory}
\label{sec:discussion-haplotype}

Haplotype blocks~\cite{ebert2021haplotype} are regions of the genome where
linkage disequilibrium is high and recombination is rare.  Within a haplotype
block, the set of alleles at variant sites is highly correlated, and blocks
are inherited as units across generations.

Our coalitions are, in a sense, reference-free approximations of haplotype
blocks.  A coalition represents a set of k-mer modules that co-occur on a
single read, which (if the read spans a haplotype block) captures the allelic
composition of that block.  The key difference is resolution: a haplotype
block may span thousands of bases, while a sequencing read spans only
$\sim$150\,bp.  Our coalitions capture only the local module structure within
a single read, not the full extent of a haplotype block.

This resolution limitation is fundamental to short-read reference-free
analysis.  Long-read sequencing (PacBio HiFi, Oxford Nanopore) would produce
reads spanning multiple kilobases, potentially capturing entire haplotype
blocks within single reads.  Coalition Transfer on long-read data could yield
dramatically improved results, as each coalition would represent a much larger
and more informative genomic unit.

\subsection{The Symbiogenesis Contribution: Is Coalition Transfer Genuinely Different?}
\label{sec:discussion-symbiogenesis}

A critical question is whether the symbiogenesis framing---modules as organisms,
coalitions as symbiotic communities---provides genuine algorithmic insight or is
merely an evocative metaphor.

The honest answer is mixed.  On one hand, the symbiogenesis framing motivated
the development of Coalition Transfer: the idea that ``organisms cooperate in
coalitions'' led directly to the idea that ``modules cooperate in read-level
combinations''.  On the other hand, the algorithm itself is a straightforward
application of standard techniques: co-occurrence counting, connected-component
discovery, set intersection, and IDF weighting.  None of these require the
symbiogenesis metaphor; they could have been derived from first principles in
information retrieval or combinatorics.

The metaphor's value may lie more in \emph{direction} than in \emph{mechanism}.
It suggests specific research directions---boundary conservation, ecological
succession of coalitions across generations, niche differentiation of module
types---that might not arise from a purely information-theoretic framing.
Whether these directions yield practical improvements remains to be seen.

\subsection{Computational Advantages}
\label{sec:discussion-computation}

Coalition Transfer's most unambiguous advantage over Module Persistence~v2 is
computational efficiency.  At 272 seconds total (44 seconds per sample),
Coalition Transfer is $4{\times}$ faster than Module Persistence~v2 (1{,}085
seconds, 180 seconds per sample).  The memory improvement is even more dramatic:
approximately 0\,MB heap delta versus 4{,}806\,MB for Module Persistence~v2.

The speed advantage comes from three sources:

\begin{enumerate}
  \item \textbf{Larger windows} reduce the number of co-occurrence updates
        during module extraction.
  \item \textbf{No graph eigendecomposition.}  Module Persistence~v2 computes
        spectral features of the module co-occurrence graph, which requires
        matrix operations.  Coalition Transfer replaces this with the simpler
        integrity cosine.
  \item \textbf{Efficient profiling.}  Coalition profiling is a single-pass
        counting operation over pre-extracted module assignments, requiring
        only a hash map from fingerprints to counts.
\end{enumerate}

The memory advantage arises because Coalition Transfer does not need to
materialize the full module co-occurrence graph.  Module Persistence~v2 builds
a dense adjacency matrix for spectral analysis; Coalition Transfer discards
the graph after connected-component extraction and operates on the much more
compact coalition profiles.

For practical deployment---especially in resource-constrained environments or
on consumer hardware---this efficiency advantage is significant.  A 4.5-minute
pipeline that uses negligible memory is feasible on a laptop; an 18-minute
pipeline requiring 5\,GB of RAM is not.

\subsection{Limitations}
\label{sec:discussion-limitations}

Coalition Transfer has several significant limitations:

\begin{enumerate}
  \item \textbf{Negative weakest gap.}  The in-law pair
        Mat.GF~$\leftrightarrow$~Father scores above the weakest parent--child
        pair, precluding threshold-based classification.  This is the most
        critical limitation for any practical application.

  \item \textbf{Non-discriminative integrity.}  The integrity cosine contributes
        nothing to separation because all values are mechanically~1.0.  This
        wastes 30\% of the scoring function's capacity.

  \item \textbf{Exact coalition matching.}  The fingerprinting approach requires
        exact module-set identity for a coalition match.  Any variation in
        module composition creates a new coalition, reducing sensitivity to
        genuine inheritance events where boundary effects alter the module set.

  \item \textbf{Short-read limitation.}  With 150\,bp reads, each coalition
        captures only a small genomic window.  The combinatorial richness of
        coalitions is limited by the physical read length, which constrains
        the number of modules that can co-occur on a single read.

  \item \textbf{Single pedigree evaluation.}  All results are from a single
        six-member pedigree.  The method's performance on larger pedigrees,
        different populations, or more distant relationships is unknown.

  \item \textbf{Coverage sensitivity.}  The number of discovered modules
        varies 1.7-fold across samples (507--865), suggesting that coverage
        differences influence module extraction.  This variation propagates
        to the coalition level and may bias pairwise comparisons.

  \item \textbf{No population-level calibration.}  Unlike classical IBD
        methods that calibrate against a population panel, Coalition Transfer
        has no mechanism to distinguish population-level sharing from
        family-level sharing.

  \item \textbf{Parameter sensitivity.}  The method uses several parameters
        (window size, motif $k$, minimum support, minimum cohesion, scoring
        weights) that were chosen by experimentation rather than principled
        optimization.  Different parameter settings may yield substantially
        different results.
\end{enumerate}


% ════════════════════════════════════════════════════════════════════
\section{Future Work}
\label{sec:future}

The Coalition Transfer results suggest several promising research directions.

\subsection{Boundary Conservation Analysis}

The most immediate improvement would be measuring how coalition boundaries
are conserved across related individuals.  If a parent and child share a
haplotype block, the coalition boundaries at the edges of that block should
be similar (but not necessarily identical) in both samples.  A
\emph{boundary conservation score} could provide a fourth scoring component
that captures information currently lost by exact fingerprint matching.

\subsection{Ecological Role Assignment}

In ecology, species within a coalition play different roles: keystone species,
facilitators, competitors~\cite{bray1957}.  We could assign analogous roles
to modules within a coalition---e.g., a ``core module'' that appears in every
instance of the coalition vs.\ a ``peripheral module'' that appears only in
some instances.  This role structure could provide additional discriminative
features for relatedness scoring.

\subsection{Graph Diffusion over Coalition Networks}

Rather than comparing coalition profiles directly, we could construct a
\emph{coalition graph} where nodes are coalitions and edges connect
coalitions that share modules.  Random-walk-based similarity on this graph
(e.g., personalized PageRank) could capture higher-order relationships
between coalition profiles that are invisible to direct comparison.

\subsection{Multi-Resolution Coalition Analysis}

Currently, all coalitions are extracted at a single read length ($\sim$150\,bp).
A multi-resolution approach would extract coalitions at multiple window sizes
(50\,bp, 100\,bp, 150\,bp, 300\,bp) and combine the signals.  Smaller windows
would capture fine-grained local structure; larger windows (via paired-end
information or read merging) would capture broader haplotype blocks.

\subsection{GPU Acceleration}

The module extraction stage (97.4\% of total runtime) is embarrassingly
parallel: each read can be processed independently, and co-occurrence
counting can be parallelized across windows.  The yalumba GPU pipeline
(TypeScript DSL $\to$ C codegen $\to$ SPIR-V $\to$ GPU dispatch) is designed
for exactly this kind of workload.  A GPU-accelerated module extraction
kernel could reduce the 265-second extraction phase to seconds, making
interactive experimentation feasible.

\subsection{Approximate Coalition Matching}

Replacing exact fingerprint matching with locality-sensitive hashing (LSH)
would allow fuzzy coalition comparison.  Two coalitions with Jaccard overlap
$> 0.8$ of their module sets could be treated as variants of the same
underlying haplotype structure.  This would address the fundamental
specificity problem identified in Section~\ref{sec:discussion-unit} and
potentially recover much of the signal lost to boundary effects.

\subsection{Long-Read Validation}

Testing Coalition Transfer on PacBio HiFi or Oxford Nanopore reads
(10{,}000--20{,}000\,bp) would dramatically increase the number of modules
per coalition and the specificity of each coalition fingerprint.  Long reads
could capture entire haplotype blocks, making coalitions a direct analogue
of IBD segments rather than a crude approximation.


% ════════════════════════════════════════════════════════════════════
\section{Conclusion}
\label{sec:conclusion}

Coalition Transfer is the second algorithm in the yalumba symbiogenesis family,
extending Module Persistence by shifting the unit of analysis from individual
k-mer modules to composite coalitions---sets of modules that co-occur on single
sequencing reads.  The method is motivated by a biologically grounded intuition:
inherited DNA transmits not just individual genomic features but specific
\emph{combinations} of features, and these combinations should be more
informative than their individual components.

Evaluated on the CEPH~1463 pedigree (6 members, 10 pairwise comparisons),
Coalition Transfer achieves a mean separation of $+1.36\%$ between
parent--child and unrelated/in-law pairs, with a weakest gap of $-0.55\%$.
These results confirm that coalition-level sharing is correlated with
biological relatedness but fail to achieve the separation needed for
threshold-based classification.  The in-law confound
(Mat.GF~$\leftrightarrow$~Father scoring above Pat.GM~$\to$~Father) persists
from all previous experiments.

The algorithm's computational profile is favorable: 272 seconds total
($4{\times}$ faster than Module Persistence~v2) with negligible memory
overhead.  However, this speed advantage comes at the cost of reduced
separation ($+1.36\%$ vs.\ $+4.76\%$), and both symbiogenesis methods remain
far behind the Rare-Run P90 baseline ($+583\%$).

Three specific findings inform future development:

\begin{enumerate}
  \item \textbf{Integrity is non-discriminative} because exact fingerprint
        matching mechanically produces integrity~1.0.  Approximate coalition
        matching is needed to make this component meaningful.
  \item \textbf{Coalitions are too specific.}  Exact module-set matching
        misses genuine inheritance events where boundary effects alter the
        coalition composition.  Fuzzy matching (via LSH or Hamming thresholds)
        could recover lost signal.
  \item \textbf{Rarity remains the dominant signal.}  Across all yalumba
        experiments, rare features---whether k-mers, modules, or
        coalitions---provide the strongest discriminative power for relatedness
        detection.  Future algorithms should prioritize rarity-aware scoring.
\end{enumerate}

The work establishes coalitions as a meaningful unit of reference-free
genomic analysis, documents their current limitations honestly, and
identifies specific algorithmic improvements that could close the gap
with simpler but more effective baselines.  The symbiogenesis research
program continues: the metaphor of biological merger provides direction,
even when the current instantiation falls short of its promise.


% ════════════════════════════════════════════════════════════════════
\section*{References}
\label{sec:references}

\begingroup
\renewcommand{\section}[2]{}%
\bibliographystyle{unsrt}
\bibliography{../references}
\endgroup

\end{document}
